\documentclass[12pt, a4paper]{article}

% ============================================================
%  Paquetes
% ============================================================
\usepackage[utf8]{inputenc}
\usepackage[T1]{fontenc}
\usepackage[spanish]{babel}
\usepackage{amsmath, amssymb}
\usepackage[ruled, vlined, linesnumbered, spanish]{algorithm2e}
\usepackage{geometry}
\usepackage{hyperref}

\geometry{margin=2.5cm}

% ============================================================
%  Metadatos del documento
% ============================================================
\title{Colonia de Abejas Artificiales Simple (ABC) --- Pseudocódigo}
\author{}
\date{}

% ============================================================
\begin{document}
\maketitle

% ------------------------------------------------------------
\section{Introducción y fundamento biológico}
% ------------------------------------------------------------

El algoritmo \textbf{Artificial Bee Colony (ABC)} de Karaboga (2005) modela el
comportamiento de una colonia de abejas melíferas durante la búsqueda y
explotación de fuentes de alimento. En el contexto de optimización combinatoria,
cada \emph{fuente de alimento} es una solución candidata y la calidad de la
fuente es la aptitud (\emph{fitness}) de la solución.

La implementación \texttt{busqueda\_abejas\_simple} sigue de forma fiel la
estructura de Karaboga (2005) y la adapta al CARP con las siguientes
modificaciones justificadas:

\begin{enumerate}
    \item \textbf{Sesgo inter/intra-ruta}: se usa el helper compartido
          \texttt{seleccionar\_grupo\_operadores\_inter\_intra} en fases
          empleadas y observadoras para reparar violaciones de capacidad
          —sin él, el algoritmo podría quedarse atascado en regiones infactibles.
    \item \textbf{$p_{\mathrm{inter}}$ dinámico}: cuando hay violación de
          capacidad, la probabilidad de elegir operadores inter-ruta sube
          automáticamente a $\max(p_{\mathrm{inter}}, 0.8)$.
    \item \textbf{Penalización $\lambda$}: $f(s) = c(s) + \lambda \cdot
          \mathrm{viol}(s)$; permite trabajar con soluciones infactibles y guía
          la búsqueda hacia la región factible.
    \item \textbf{Criterio de parada por estancamiento}: \texttt{max\_iter\_sin\_mejora}
          (calibrado a $\max(50,\,3n)$ por instancia).
    \item \textbf{Scouts aleatorias puras}: Karaboga estricto; las fuentes
          agotadas se sustituyen por soluciones generadas al azar (no por vecinos
          de la mejor actual, como hace la versión extendida del proyecto).
    \item \textbf{Corrección del bug de imputación de mejora}: la función
          \texttt{registrar\_mejora} se invoca inmediatamente cuando un vecino
          supera el mejor global, dentro del mismo \texttt{if} de detección.
\end{enumerate}

% ------------------------------------------------------------
\section{Notación}
% ------------------------------------------------------------

\begin{itemize}
    \item $G = (V, E, R)$: grafo de la instancia CARP. $R \subseteq E$ son las
          aristas requeridas.
    \item $F = \{s_1, s_2, \ldots, s_{N}\}$: conjunto de fuentes de alimento
          (soluciones activas), con $N = \mathit{num\_fuentes}$.
    \item $s_i$: fuente $i$ (una solución CARP --- conjunto de rutas sobre $R$).
    \item $s^{*}$: mejor solución encontrada durante toda la corrida.
    \item $c(s)$: costo puro de la solución $s$ (suma de costos de travesía y
          servicio, sin penalización).
    \item $\mathrm{viol}(s)$: exceso total de demanda respecto a la capacidad del
          vehículo.
    \item $f(s) = c(s) + \lambda \cdot \mathrm{viol}(s)$: objetivo penalizado
          ($\lambda > 0$).
    \item $\mathit{trials}[i]$: contador de intentos fallidos consecutivos de
          la fuente $i$ (se reinicia a 0 cuando la fuente mejora).
    \item $L_{\mathrm{aban}}$: límite de abandono — si $\mathit{trials}[i] \geq
          L_{\mathrm{aban}}$, la fuente $i$ se reemplaza por un scout.
    \item $\mathcal{O}_{\mathrm{intra}} = \{\texttt{relocate\_intra},\,
          \texttt{swap\_intra},\, \texttt{2opt\_intra}\}$: operadores
          intra-ruta.
    \item $\mathcal{O}_{\mathrm{inter}} = \{\texttt{relocate\_inter},\,
          \texttt{swap\_inter},\, \texttt{2opt\_star},\,
          \texttt{cross\_exchange},\, \texttt{or\_opt\_2},\,
          \texttt{or\_opt\_3}\}$: operadores inter-ruta.
    \item $p_{\mathrm{inter}} \in [0,1]$: probabilidad base de elegir el grupo
          inter-ruta (cuando la solución es factible).
    \item $p_{\mathrm{efec}}$: probabilidad efectiva de elegir inter-ruta tras
          aplicar el piso dinámico.
    \item $n = |R|$: número de tareas requeridas (variable de escala
          \emph{instance-aware}).
\end{itemize}

\medskip
\noindent
\textbf{Sesgo inter/intra dinámico.} En cada ciclo se calcula la violación media
$\bar{v} = \tfrac{1}{N}\sum_{i=1}^{N}\mathrm{viol}(s_i)$. La probabilidad
efectiva de elegir el grupo inter-ruta es:
\[
p_{\mathrm{efec}} \;=\;
\begin{cases}
  \max(p_{\mathrm{inter}},\; 0.8) & \text{si } \bar{v} > 0, \\
  p_{\mathrm{inter}}              & \text{si } \bar{v} = 0.
\end{cases}
\]
Con probabilidad $p_{\mathrm{efec}}$ se selecciona $\mathcal{O}_{\mathrm{inter}}$;
en caso contrario, $\mathcal{O}_{\mathrm{intra}}$. Dentro del grupo elegido, el
operador concreto se muestrea de manera uniforme.

% ------------------------------------------------------------
\section{Calibración \emph{instance-aware} de parámetros}
% ------------------------------------------------------------

Los parámetros del algoritmo se escalan con $n = |R|$ para evitar que
corridas demasiado cortas (instancias pequeñas) o demasiado largas (instancias
grandes) distorsionen la comparación. Las fórmulas por defecto son:

\begin{align*}
\mathit{iteraciones}        &= \max(200,\; 20 \cdot n) \\
\mathit{num\_fuentes}       &= \max(10,\; \lceil 2\sqrt{n} \rceil) \\
\mathit{limite\_abandono}   &= \max(15,\; \lfloor n/2 \rfloor) \\
\mathit{max\_iter\_sin\_mejora} &= \max(50,\; 3 \cdot n)
\end{align*}

El usuario puede sobrescribir cualquier parámetro con un \emph{factor de escala}
$f$ (e.g.\ $\mathit{factor\_iter} = 20$, lo que equivale a
$\mathit{iteraciones} = \max(200, 20n)$) o con un valor absoluto. La precedencia es:

\[
\text{valor absoluto} \;\succ\; \text{factor de escala} \;\succ\; \text{fórmula por defecto}
\]

% ------------------------------------------------------------
\section{Pseudocódigos}
% ------------------------------------------------------------

% -------- Algoritmo ABC Simple completo ---------------------
\begin{algorithm}[H]
\DontPrintSemicolon
\caption{ABC Simple para CARP (Karaboga 2005 adaptado)}
\label{alg:abc_simple}

\KwIn{Instancia CARP $G$;\; solución inicial $s_0$;\;
      $p_{\mathrm{inter}}$;\; penalización $\lambda$;\;
      opcionalmente: $N$, $L_{\mathrm{aban}}$, $\mathit{iteraciones}$,
      $\mathit{max\_iter\_sin\_mejora}$ (todos instance-aware si se omiten)}
\KwOut{Mejor solución encontrada $s^{*}$}

\BlankLine
\tcp{Calibración instance-aware (valores default si no se proporcionan)}
$n \leftarrow |R|$ \tcp{número de tareas requeridas}
$\mathit{iteraciones} \leftarrow \max(200,\; 20 \cdot n)$ \;
$N \leftarrow \max(10,\; \lceil 2\sqrt{n} \rceil)$ \;
$L_{\mathrm{aban}} \leftarrow \max(15,\; \lfloor n/2 \rfloor)$ \;
$M \leftarrow \max(50,\; 3 \cdot n)$
    \tcp{máx. iteraciones sin mejora}

\BlankLine
\tcp{Inicialización de fuentes}
$\mathit{fuentes}[0] \leftarrow s_0$;\quad
$\mathit{trials}[0] \leftarrow 0$ \;
\For{$i \leftarrow 1$ \KwTo $N - 1$}{
    $\mathit{fuentes}[i] \leftarrow$ \textsc{GenerarVecino}($\mathit{fuentes}[0],\,
        \mathcal{O}_{\mathrm{inter}} \cup \mathcal{O}_{\mathrm{intra}}$) \;
    $\mathit{trials}[i] \leftarrow 0$ \;
}

\BlankLine
\tcp{Mejor inicial (mínimo objetivo penalizado sobre todas las fuentes)}
$s^{*} \leftarrow \arg\min_{i} f(\mathit{fuentes}[i])$ \;
$f^{*} \leftarrow f(s^{*})$ \;
$\mathit{iter\_sin\_mejora} \leftarrow 0$ \;

\BlankLine
\tcp{Precómputo UNA VEZ de las particiones intra/inter (microoptimización)}
Construir $\mathcal{O}_{\mathrm{intra}}$, $\mathcal{O}_{\mathrm{inter}}$,
$\mathcal{O}_{\mathrm{fallback}}$ a partir de los operadores activos \;

\BlankLine
\For{$\mathit{ciclo} \leftarrow 1$ \KwTo $\mathit{iteraciones}$}{

    \tcp{Criterio de parada anticipada por estancamiento}
    \If{$\mathit{iter\_sin\_mejora} \geq M$}{
        \textbf{break} \;
    }

    $f^{*}_{\mathrm{inicio}} \leftarrow f^{*}$ \tcp{snapshot del inicio del ciclo}

    \BlankLine
    \tcp{Violación media y p\_inter efectivo para este ciclo}
    $\bar{v} \leftarrow \tfrac{1}{N}\sum_{i=1}^{N} \mathrm{viol}(\mathit{fuentes}[i])$ \;
    \lIf{$\bar{v} > 0$}{$p_{\mathrm{efec}} \leftarrow \max(p_{\mathrm{inter}},\; 0.8)$}
    \lElse{$p_{\mathrm{efec}} \leftarrow p_{\mathrm{inter}}$}

    \BlankLine
    \tcp{\textbf{Fase 1 — Empleadas} (explotación individual)}
    \For{$i \leftarrow 0$ \KwTo $N - 1$}{
        $\mathcal{O}_{\mathrm{grupo}} \leftarrow$ \textsc{ElegirGrupo}($p_{\mathrm{efec}},\,
            \mathcal{O}_{\mathrm{inter}},\, \mathcal{O}_{\mathrm{intra}}$) \;
        $(s',\, m) \leftarrow$ \textsc{GenerarVecino}($\mathit{fuentes}[i],\,
            \mathcal{O}_{\mathrm{grupo}}$) \;
        \textsc{Proponer}($m$) \;
        \eIf{$f(s') < f(\mathit{fuentes}[i]) - \varepsilon$}{
            $\mathit{fuentes}[i] \leftarrow s'$;\quad
            $\mathit{trials}[i] \leftarrow 0$ \;
            \textsc{Aceptar}($m$) \;
            \tcp{Detección inmediata de mejora global (corrección del bug de imputación)}
            \If{$f(s') < f^{*} - \varepsilon$}{
                $s^{*} \leftarrow s'$;\quad
                $f^{*} \leftarrow f(s')$ \;
                \textsc{RegistrarMejora}($m$) \;
            }
        }{
            $\mathit{trials}[i] \leftarrow \mathit{trials}[i] + 1$ \;
        }
    }

    \BlankLine
    \tcp{\textbf{Fase 2 — Observadoras} (explotación ponderada por ruleta)}
    \tcp{Recalcular violación media y $p_{\mathrm{efec}}$ tras las empleadas}
    $\bar{v} \leftarrow \tfrac{1}{N}\sum_{i=1}^{N} \mathrm{viol}(\mathit{fuentes}[i])$ \;
    \lIf{$\bar{v} > 0$}{$p_{\mathrm{efec}} \leftarrow \max(p_{\mathrm{inter}},\; 0.8)$}
    \lElse{$p_{\mathrm{efec}} \leftarrow p_{\mathrm{inter}}$}

    \BlankLine
    \tcp{Probabilidades de la ruleta: fitness inverso normalizado}
    \For{$i \leftarrow 0$ \KwTo $N - 1$}{
        $w_i \leftarrow \dfrac{1}{1 + f(\mathit{fuentes}[i])}$ \;
    }
    $W \leftarrow \sum_{i=0}^{N-1} w_i$;\quad
    $P_i \leftarrow w_i / W$ para todo $i$ \;

    \BlankLine
    \tcp{Muestrear $N$ índices CON REEMPLAZO (fuentes buenas reciben más visitas)}
    $\mathit{idxs} \leftarrow \textsc{RuletaMuestreaConReemplazo}(P, k=N)$ \;

    \BlankLine
    \tcp{Generar lote de vecinos (uno por cada índice muestreado)}
    \For{$k \leftarrow 0$ \KwTo $N - 1$}{
        $i \leftarrow \mathit{idxs}[k]$ \;
        $\mathcal{O}_{\mathrm{grupo}} \leftarrow$ \textsc{ElegirGrupo}($p_{\mathrm{efec}},\,
            \mathcal{O}_{\mathrm{inter}},\, \mathcal{O}_{\mathrm{intra}}$) \;
        $(s'_k,\, m_k) \leftarrow$ \textsc{GenerarVecino}($\mathit{fuentes}[i],\,
            \mathcal{O}_{\mathrm{grupo}}$) \;
    }

    \BlankLine
    \tcp{Evaluación vectorizada del lote (GPU si disponible)}
    $\{f(s'_k)\}_{k=0}^{N-1} \leftarrow$ \textsc{EvaluarLote}($\{s'_k\}$)
        \tcp{\texttt{costo\_lote\_penalizado\_ids}}

    \BlankLine
    \tcp{Actualización greedy de fuentes}
    \For{$k \leftarrow 0$ \KwTo $N - 1$}{
        $i \leftarrow \mathit{idxs}[k]$ \;
        \textsc{Proponer}($m_k$) \;
        \eIf{$f(s'_k) < f(\mathit{fuentes}[i]) - \varepsilon$}{
            $\mathit{fuentes}[i] \leftarrow s'_k$;\quad
            $\mathit{trials}[i] \leftarrow 0$ \;
            \textsc{Aceptar}($m_k$) \;
            \tcp{Detección inmediata de mejora global}
            \If{$f(s'_k) < f^{*} - \varepsilon$}{
                $s^{*} \leftarrow s'_k$;\quad $f^{*} \leftarrow f(s'_k)$ \;
                \textsc{RegistrarMejora}($m_k$) \;
            }
        }{
            $\mathit{trials}[i] \leftarrow \mathit{trials}[i] + 1$ \;
        }
    }

    \BlankLine
    \tcp{\textbf{Fase 3 — Scouts} (diversificación con soluciones aleatorias puras)}
    \ForEach{$i$ \textbf{ tal que } $\mathit{trials}[i] \geq L_{\mathrm{aban}}$}{
        $s_{\mathrm{nuevo}} \leftarrow$ \textsc{SoluciónAleatoria}($G$)
            \tcp{barajar tareas + asignación greedy de capacidad}
        $\mathit{fuentes}[i] \leftarrow s_{\mathrm{nuevo}}$ \;
        $\mathit{trials}[i] \leftarrow 0$ \;
        \tcp{El scout NO llama a Proponer/Aceptar: no hay operador de vecindad}
        \tcp{Si el scout mejoró el global, se registra sin imputar operador}
        \If{$f(s_{\mathrm{nuevo}}) < f^{*} - \varepsilon$}{
            $s^{*} \leftarrow s_{\mathrm{nuevo}}$;\quad
            $f^{*} \leftarrow f(s_{\mathrm{nuevo}})$ \;
            \tcp{Mejora contabilizada pero sin asociar a ningún operador}
        }
    }

    \BlankLine
    \tcp{Actualización del contador de estancamiento}
    \lIf{$f^{*} < f^{*}_{\mathrm{inicio}} - \varepsilon$}{%
        $\mathit{iter\_sin\_mejora} \leftarrow 0$}
    \lElse{%
        $\mathit{iter\_sin\_mejora} \leftarrow \mathit{iter\_sin\_mejora} + 1$}
}

\BlankLine
\Return $s^{*}$
\end{algorithm}

\bigskip

% -------- Auxiliar: Solución aleatoria (scouts canónicos) ----
\begin{algorithm}[H]
\DontPrintSemicolon
\caption{Generación de solución aleatoria para scouts}
\label{alg:solucion_aleatoria}

\KwIn{Instancia CARP $G$;\; capacidad del vehículo $Q$;\; depósito $d$}
\KwOut{Solución CARP aleatoria $s$}

\BlankLine
$\mathit{tareas} \leftarrow$ barajar aleatoriamente $R$ \;
$s \leftarrow [\,]$;\quad $\mathit{ruta\_actual} \leftarrow [d]$;\quad
$\mathit{carga\_actual} \leftarrow 0$ \;
\ForEach{tarea $t \in \mathit{tareas}$}{
    \eIf{$\mathit{carga\_actual} + \mathrm{demanda}(t) \leq Q$}{
        $\mathit{ruta\_actual} \leftarrow \mathit{ruta\_actual} + [t]$ \;
        $\mathit{carga\_actual} \leftarrow \mathit{carga\_actual} + \mathrm{demanda}(t)$ \;
    }{
        $s \leftarrow s + [\mathit{ruta\_actual} + [d]]$
            \tcp{cerrar ruta actual con retorno al depósito}
        $\mathit{ruta\_actual} \leftarrow [d, t]$ \;
        $\mathit{carga\_actual} \leftarrow \mathrm{demanda}(t)$ \;
    }
}
\lIf{$|\mathit{ruta\_actual}| > 1$}{
    $s \leftarrow s + [\mathit{ruta\_actual} + [d]]$}
\Return $s$
\end{algorithm}

% ------------------------------------------------------------
\section{Notas de implementación}
% ------------------------------------------------------------

\paragraph{Evaluación vectorizada (GPU).}
En la fase observadoras, los $N$ vecinos generados en lote se evalúan en una
sola llamada a \texttt{costo\_lote\_penalizado\_ids}, que empaqueta los arrays
de distancias y demandas en NumPy y, si hay GPU disponible (\texttt{usar\_gpu=True}),
delega la reducción a CuPy. En la fase empleadas cada vecino se evalúa
individualmente con \texttt{costo\_rapido\_ids} porque la decisión de
aceptar/rechazar se toma por fuente (la vectorización no daría speedup allí).

\paragraph{Corrección del bug de imputación.}
En la versión extendida \texttt{busqueda\_abejas}, la detección de mejora global
se hacía comparando $f(s)$ antes y después del ciclo completo, y la mejora se
atribuía al \emph{último movimiento aceptado} en el ciclo (que podría ser
cualquier movimiento posterior al que realmente mejoró el mejor). Aquí, la
llamada a \textsc{RegistrarMejora}$(m)$ ocurre dentro del mismo \texttt{if}
que actualiza $s^{*}$, garantizando que $\sum_{\mathrm{op}}\mathit{mejoraron}[\mathrm{op}]
\leq \mathit{mejoras}$, con la diferencia exacta contabilizada a los scouts.

\paragraph{Ruleta de observadoras.}
El fitness de la fuente $i$ en la fórmula de Karaboga es $\mathit{fit}_i = 1/(1 + f_i)$
(válida para $f_i \geq 0$, siempre cumplida en CARP). Las probabilidades de
selección son $P_i = \mathit{fit}_i / \sum_j \mathit{fit}_j$. La selección se
realiza con Python's \texttt{random.choices} con repetición ($k = N$), lo que
permite que las mejores fuentes sean visitadas más de una vez por ciclo.

\paragraph{Inicialización de fuentes.}
La fuente 0 se siembra con la solución inicial de referencia. Las fuentes
$1 \ldots N-1$ se siembran con vecinos aleatorios de la fuente 0 (en lugar de
soluciones completamente aleatorias) para arrancar con diversidad pero cercana
a una buena solución de partida. Es el compromiso estándar de la literatura
aplicada de ABC.

\end{document}
